\subsection{积分属于 E 的判据}

命题 7. --- 设 E 是一个 Hausdorff 局部凸空间，且 \(\mathbf{f} \in \mathcal{K}(\mathrm{X};\mathrm{E})\) 。若 \(\mathbf{f}(\mathrm{X})\) 含于 E 的一个完备凸子集 A，则 \(\int \mathbf{f} d\mu \in \mathbf{E}\) 。

设 K 是 f 的支集，按假设它是紧的。由于 f 在 X - K 上为零， \(\mathbf{f}(X)\) 等于 \(\mathbf{f}(K)\) 或 \(\mathbf{f}(K) \cup \{0\}\) ，因 f 连续且 E 是 Hausdorff 的，故它是紧的；于是 \(\mathbf{f}(X)\) 在 E 中的闭凸包 C 是预紧的（对由 E 的一致结构诱导的一致结构而言）（TVS, II, §4, No.~1, 命题 3）。但由于 C 是完备空间 A 的一个闭子集，故 C 是完备的从而是紧的；更有甚者，C 对弱化拓扑 \(\sigma(\mathrm{E}, \mathrm{E}')\) 是紧的；由于该拓扑是由 \(\sigma(\mathrm{E}'*, \mathrm{E}')\) 诱导的，故 C 也是 \(\mathbf{f}(X)\) 在 \(E'\) 中对拓扑 \(\sigma(\mathrm{E}'*, \mathrm{E}')\) 的闭凸包；于是证明由 No.~2 命题 4 的推论完成。

推论 1. --- 设 E 是一个 Hausdorff 局部凸空间；对每个函数 \(\mathbf{f} \in \mathcal{K}(\mathrm{X};\mathrm{E})\) ， \(\int \mathbf{f} d\mu\) 属于 E 的完备化 \(\widehat{\mathbf{E}}\) 。

由于 E 与 \(\widehat{E}\) 的对偶相同，只需把 f 看作取值于 \(\widehat{E}\) 而应用命题 7。

推论 2. --- 若 E 是一个拟完备的 Hausdorff 局部凸空间，则 \(\int \mathbf{f} d\mu \in \mathrm{E}\) 对每个函数 \(\mathbf{f} \in \mathcal{K}(\mathrm{X};\mathrm{E})\) 成立。

如命题 7 证明开头所指出的， \(\mathbf{f}(\mathrm{X})\) 是紧的，且其在 E 中的闭凸包 C 是预紧的，从而是有界的；但由于集 C 是闭的且有界，按假设它是完备的，于是只需应用命题 7。

在 Ch. VI, §1, No.~2 中，我们将看到 \(\int \mathbf{f} d\mu\) 属于 E 的其他判据，它们特别适用于 \(\widetilde{\mathcal{K}}(\mathrm{X};\mathrm{E})\) 中的函数，而不仅限于 \(\mathcal{K}(\mathrm{X};\mathrm{E})\) 中的函数。

命题 7 的推论 2 可应用于下述两种情形：\(1^{\circ}\) E 是 Banach 空间；\(2^{\circ}\) E 是某个桶式 Hausdorff 局部凸空间 G 的对偶，且 E 赋以某个 S-拓扑，其中 S 是 G 由有界子集组成的一个覆盖（TVS, III, §4, No.~2, Th. 1 的推论 4）。例如，当 E 是 Banach 空间的弱对偶，或是赋以淡拓扑的测度空间 \(\mathcal{M}(\mathrm{Y};\mathbf{C})\) 时，可以应用命题 7 的推论 2。

若 \(\mathbf{X} = \mathbf{R}\) ， \(\mu\) 是 \(\mathbf{R}\) 上的 Lebesgue 测度，且 \(\mathrm{E}\) 是 Banach 空间，则积分 \(\int \mathbf{f} d\mu\) （ \(\mathcal{K}(\mathrm{X};\mathrm{E})\) 中函数的）不是别的，正是在

\[
\int_ {- \infty} ^ {+ \infty} \mathbf {f} (x) d x
\]

中定义的积分，即 FRV, II, §2, No.~1 中的积分；这由公式 (1) 及 FRV, I, §1, No.~2, 命题 2 的推论得出。

